%%% The main file. It contains definitions of basic parameters and includes all other parts.

% Meta-data of your thesis (please edit)
\input metadata.tex

% Generate metadata in XMP format for use by the pdfx package
\input xmp.tex

%% Settings for single-side (simplex) printing
% Margins: left 40mm, right 25mm, top and bottom 25mm
% (but beware, LaTeX adds 1in implicitly)
\documentclass[12pt,a4paper]{report}
\setlength\textwidth{145mm}
\setlength\textheight{247mm}
\setlength\oddsidemargin{15mm}
\setlength\evensidemargin{15mm}
\setlength\topmargin{0mm}
\setlength\headsep{0mm}
\setlength\headheight{0mm}
% \openright makes the following text appear on a right-hand page
\let\openright=\clearpage

%% Settings for two-sided (duplex) printing
% \documentclass[12pt,a4paper,twoside,openright]{report}
% \setlength\textwidth{145mm}
% \setlength\textheight{247mm}
% \setlength\oddsidemargin{14.2mm}
% \setlength\evensidemargin{0mm}
% \setlength\topmargin{0mm}
% \setlength\headsep{0mm}
% \setlength\headheight{0mm}
% \let\openright=\cleardoublepage

%% If the thesis has no printed version, symmetric margins look better
% \documentclass[12pt,a4paper]{report}
% \setlength\textwidth{145mm}
% \setlength\textheight{247mm}
% \setlength\oddsidemargin{10mm}
% \setlength\evensidemargin{10mm}
% \setlength\topmargin{0mm}
% \setlength\headsep{0mm}
% \setlength\headheight{0mm}
% \let\openright=\clearpage

%% Generate PDF/A-2u
\usepackage[a-2u]{pdfx}

%% Prefer Latin Modern fonts
\usepackage{lmodern}
% If we are not using LuaTeX, we need to set up character encoding:
\usepackage{iftex}
\ifpdftex
\usepackage[utf8]{inputenc}
\usepackage[T1]{fontenc}
\usepackage{textcomp}

\fi

%% Further useful packages (included in most LaTeX distributions)
\usepackage{amsmath}        % extensions for typesetting of math
\usepackage{amsfonts}       % math fonts
\usepackage{amsthm}         % theorems, definitions, etc.
\usepackage{bm}             % boldface symbols (\bm)
\usepackage{booktabs}       % improved horizontal lines in tables
\usepackage{caption}        % custom captions of floating objects
\usepackage{dcolumn}        % improved alignment of table columns
\usepackage{floatrow}       % custom float environments
\usepackage{graphicx}       % embedding of pictures
\usepackage{indentfirst}    % indent the first paragraph of a chapter
\usepackage[nopatch=item]{microtype}   % micro-typographic refinement
\usepackage{paralist}       % improved enumerate and itemize
\usepackage[nottoc]{tocbibind} % makes sure that bibliography and the lists
			    % of figures/tables are included in the table
			    % of contents
\usepackage{xcolor}         % typesetting in color

% The hyperref package for clickable links in PDF and also for storing
% metadata to PDF (including the table of contents).
% Most settings are pre-set by the pdfx package.
\hypersetup{unicode}
\hypersetup{breaklinks=true}

% Packages for computer science theses
\usepackage{algpseudocode}  % part of algorithmicx package
\usepackage{algorithm}
\usepackage{fancyvrb}       % improved verbatim environment
\usepackage{listings}       % pretty-printer of source code
\usepackage{dirtytalk}
% You might want to use cleveref for references
% \usepackage{cleveref}

% Set up formatting of bibliography (references to literature)
% Details can be adjusted in macros.tex.
%
% BEWARE: Different fields of research and different university departments
% have their own customs regarding bibliography. Consult the bibliography
% format with your supervisor.
%
% The basic format according to the ISO 690 standard with numbered references
\usepackage[natbib,style=iso-numeric,sorting=none]{biblatex}
% ISO 690 with alphanumeric references (abbreviations of authors' names)
%\usepackage[natbib,style=iso-alphabetic]{biblatex}
% ISO 690 with references Author (year)
%\usepackage[natbib,style=iso-authoryear]{biblatex}
%
% Some fields of research prefer a simple format with numbered references
% (sorting=none tells that bibliography should be listed in citation order)
%\usepackage[natbib,style=numeric,sorting=none]{biblatex}
% Numbered references, but [1,2,3,4,5] is compressed to [1-5]
%\usepackage[natbib,style=numeric-comp,sorting=none]{biblatex}
% A simple format with alphanumeric references:
%\usepackage[natbib,style=alphabetic]{biblatex}

% Load the file with bibliography entries
\addbibresource{bibliography.bib}

% Our definitions of macros (see description inside)
\input macros.tex

%%% Title page and various mandatory informational pages
\begin{document}
%%% Title page of the thesis and other mandatory pages

%%% Inscriptions at the opening page of the hard cover

% We usually do not typeset the hard cover, but if you want to do it, change \iffalse to \iftrue
\iffalse

\pagestyle{empty}
\hypersetup{pageanchor=false}
\begin{center}

\large
Charles University

\medskip

Faculty of Mathematics and Physics

\vfill

{\huge\bf\ThesisTypeTitle}

\vfill

{\huge\bf\ThesisTitle\par}

\vfill
\vfill

\hbox to \hsize{\YearSubmitted\hfil \ThesisAuthor}

\end{center}

\newpage\openright
\setcounter{page}{1}

\fi

%%% Title page of the thesis

\pagestyle{empty}
\hypersetup{pageanchor=false}
\begin{center}

\centerline{\mbox{\includegraphics[width=166mm]{img/logo-en.pdf}}}

\vspace{-8mm}
\vfill

{\bf\Large\ThesisTypeTitle}

\vfill

{\LARGE\ThesisAuthor}

\vspace{15mm}

{\LARGE\bfseries\ThesisTitle\par}

\vfill

\Department

\vfill

{
\centerline{\vbox{\halign{\hbox to 0.45\hsize{\hfil #}&\hskip 0.5em\parbox[t]{0.45\hsize}{\raggedright #}\cr
Supervisor of the \ThesisTypeName{} thesis:&\Supervisor \cr
\ifx\ThesisType\TypeRig\else
\noalign{\vspace{2mm}}
Study programme:&\StudyProgramme \cr
\fi
}}}}

\vfill

Prague \YearSubmitted

\end{center}

\newpage

%%% A page with a solemn declaration to the thesis

\openright
\hypersetup{pageanchor=true}
\vglue 0pt plus 1fill

\noindent
I declare that I carried out this \ThesisTypeName{} thesis on my own, and only with the cited
sources, literature and other professional sources.
I understand that my work relates to the rights and obligations under the Act No.~121/2000 Sb.,
the Copyright Act, as amended, in particular the fact that the Charles
University has the right to conclude a license agreement on the use of this
work as a school work pursuant to Section 60 subsection 1 of the Copyright~Act.

\vspace{10mm}

\hbox{\hbox to 0.5\hsize{%
In \hbox to 6em{\dotfill} date \hbox to 6em{\dotfill}
\hss}\hbox to 0.5\hsize{\dotfill\quad}}
\smallskip
\hbox{\hbox to 0.5\hsize{}\hbox to 0.5\hsize{\hfil Author's signature\hfil}}

\vspace{20mm}
\newpage

%%% Dedication

\openright

\noindent
\Dedication

\newpage

%%% Mandatory information page of the thesis

\openright
{\InfoPageFont

\vtop to 0.5\vsize{
\setlength\parindent{0mm}
\setlength\parskip{5mm}

Title:
\ThesisTitle

Author:
\ThesisAuthor

\DeptType:
\Department

Supervisor:
\Supervisor, \SupervisorsDepartment

Abstract:
\Abstract

Keywords:
{\def\sep{\unskip, }\ThesisKeywords}

\vfil
}

% In Czech study programmes, it is mandatory to include Czech meta-data:

\ifx\StudyLanguage\LangCS

\vtop to 0.49\vsize{
\setlength\parindent{0mm}
\setlength\parskip{5mm}

Název práce:
\ThesisTitleCS

Autor:
\ThesisAuthor

\DeptTypeCS:
\DepartmentCS

Vedoucí bakalářské práce:
\Supervisor, \SupervisorsDepartmentCS

Abstrakt:
\AbstractCS

Klíčová slova:
{\def\sep{\unskip, }\ThesisKeywordsCS}

\vfil
}

\fi

}

\newpage

%%% Further pages will be numbered
\pagestyle{plain}


%%% A page with automatically generated table of contents of the thesis

\tableofcontents

%%% Each chapter is kept in a separate file
\chapter*{Introduction}
\addcontentsline{toc}{chapter}{Introduction}
I chose the topic of my thesis because of my interest in human-robot interaction. The work focuses on gesture based control, as I believe this approach is the most natural and intuitive method of communicating with robots.\par
A mobile robot manipulator was used to test the control for performing the Pick and Place task in a real environment.\par
The topic is specified in the first chapter, where basic concepts are introduced.\par
The second chapter briefly reviews the history of computer vision, the development of gesture recognition, as well as recent work related to deictic gestures.\par
Gesture control is analyzed in the third chapter, which contains an overview of the types of gestures used and the limitations of gesture recognition with the ORBBEC Astra camera.\par
The control design is described in the next chapter. The structure of the system and its main components, the vision system and the mobile robot manipulator, are described here. Hardware information is also included.\par
The fifth chapter contains a description of the implementation process and installation instructions.
The results of the experiments are presented in the last chapter.\par
\chapter{Task Analysis}

%An~example citation: \cite{Andel07}

\section{Theoretical Background}

\subsection{Basic Definitions}
TODO: Definition of key words:
gesture recognition, object manipulation, autonomous control

\subsubsection{Gesture Recognition}
Gesture recognition technology uses a computer and a sensor to interpret human body movements. It allows direct control of machines without mechanical devices such as a keyboard or joystick.\par
Body movements are captured by a sensor. A static gesture can be identified in a single frame of raw sensor data, dynamic has to be tracked in consecutive frames during the movement.\par
The gesture is assigned to one of the predefined gesture types and the corresponding gesture type is translated into machine commands.\par
The gesture recognition process can be divided into the following parts:\par
\begin{itemize}
	\item sensor data collection
    \item gesture identification
    \item gesture tracking
    \item gesture classification
    \item gesture mapping
\end{itemize}

TODO: ref \\
Hongyi Liu, Lihui Wang (2017). Gesture recognition for human-robot collaboration: A review.\\
International Journal of Industrial Ergonomics, Volume 68, November 2018 (355-367).
https://doi.org/10.1016/j.ergon.2017.02.004\\

\subsubsection{Deictic Gesture}
A deictic gesture is a gesture that indicates direction or location from the perspective of the person performing the gesture.\par
The meaning of the deictic gesture depends on the context, it can refer to a real or a virtual environment. It could often be expressed by adverbs such as "here" and "there" or by demonstrative pronouns such as "this" and "that". We use it to specify direction and location or to identify a person or an object in the environment.\par
The most common deictic gesture is the pointing gesture. Other examples are
gestures based on head movements or eye gaze.\par

\subsubsection{Pointing Gesture}
A pointing gesture is performed by extending the arm in the appropriate
direction, usually using the index finger or hand to indicate the direction.\par
It may represent the pointing ray, which is given by, for example,
the eyes (as the origin) and the index finger, or it may have a more symbolic
meaning, such as when a person is pointing outside his field of vision or in a virtual environment.\par
Pointing with the index finger is a cross-cultural behavior. Infants most commonly use their index fingers for tactile exploration of their environment and they often use the gesture of the extended
index finger for a variety of purposes before they acquire its social meaning.\par
 
\subsection{Human-Robot Interaction (HRI)}
TODO:\par
brief description of HRI; \\

remote vs. proximate interactions; \\

roles of humans and robots in interaction: Supervisor, Operator, Mechanic, Peer, Bystander, Information consumer, Mentor (taxonomy from paper: M. A. X. Goodrich and A. C. Schultz, “Human-Robot Interaction: A Survey,” Foundations and Trends R© in Human- Computer Interaction, vol. 1, no. 3, pp. 203–275, 2007.); \\

areas of application: industrial, search and rescue, medical, social, ...\\

\section{Task description}
TODO:\par
\subsection{Pick and Place Task}

Performing 'Pick and Place' using pointing gestures:\\ \\
Pick:\\
\begin{itemize}
\item {Determine a object that was selected with a pointing gesture}
\item {Navigate close to the object}
\item {Identify the object and compute its exact coordinates}
\item {Pick the object}
\end{itemize}


Place:\\
\begin{itemize}
\item {Determine a target location from a pointing gesture}
\item {Navigate close to the location}
\item {Place the object to the location}
\end{itemize}

\subsection{Task Specification}

Requirements and restrictions:\\
gesture recognition based on image processing (available sensors - depth cameras), proximate robot control by single user, no interaction with other robots or humans, static indoor environment (robotic lab), safe manipulation with objects, safe navigation (obstacle avoidance without unnecessary emergency braking), ...

\section{Goals}
TODO:\par
\subsection{Implementation of Mobile Manipulator}
Design and implement a mobile manipulator that performs 'Pick and Place' tasks according to the given requirements;
\subsection{Comparison of deictic gestures types}
Metric: the distance between the correct coordinates (of the selected object or location) and the intersection of the pointing ray and the floor.\\ \\
Experiment with different ways of using deictic gestures:\\
- a pointing ray calculated from a pair of skeleton coordinates (head - hand, elbow - wrist, shoulder - wrist)\\
- pointing with or without visual feedback (pointed ray shown in rViz)\\



\chapter{The state of the art}

TODO: whole chapter\\
\section{History of Gesture Recognition}
Summary of gesture recognition techniques;
historical development of sensors;
...

\section{Localization and navigation with deictic gestures}

These are some (not all) examples of what I want to mention here:\\
 
Deictic gestures for multi-robot systems \\

Paper: \\
B. Gromov, L. M. Gambardella and G. A. Di Caro, "Wearable multi-modal interface for human multi-robot interaction," 2016 IEEE International Symposium on Safety, Security, and Rescue Robotics (SSRR), Lausanne, Switzerland, 2016, pp. 240-245, doi: 10.1109/SSRR.2016.7784305.\\

Use of the pointing gesture for localization \\

Paper: \\
B. Gromov, L. Gambardella, and A. Giusti. Robot Identification and Localization with Pointing Gestures. IEEE/RSJ International Conference on Intelligent Robots and Systems (IROS), 2018, pp. 3921–3928 https://people.idsia.ch/~gromov/repository/gromov2018robot.pdf \\

3D Motion planning with pointing gestures \\

Paper: \\
B. Gromov, J. Guzzi, L. Gambardella, and A. Giusti. Intuitive 3D Control of a Quadrotor in User Proximity with Pointing Gestures. IEEE International Conference on Robotics and Automation (ICRA), 2020 https://people.idsia.ch/~gromov/repository/gromov2020intuitive.pdf

\section{Interpretation of gestures}
Papers:\\

Chaudhary, A (2018). Robust Hand Gesture Recognition for Robotic Hand Control. Springer. ISBN 978-981-10-4798-5 https://doi.org/10.1007/978-981-10-4798-5\\

Alikhani, M., Khalid, B., Shome, R., Mitash, C., Bekris, K.E., Stone, M. (2020). That and There: Judging the Intent of Pointing Actions with Robotic Arms. AAAI.https://ojs.aaai.org//index.php/AAAI/article/view/6601\\

\section{Object detection with pointing gestures and speech recognition}
Li-Heng Lin, Yuchen Cui, Yilun Hao, Fei Xia, Dorsa Sadigh (2023). Gesture-Informed Robot Assistance via Foundation Models. https://arxiv.org/abs/2309.02721 \\

A. Ekrekli, A. Angleraud, G. Sharma, R. Pieters (2023). Co-speech gestures for human-robot collaboration. https://arxiv.org/abs/2311.18285 \\



% 

% 

% 
%

% paper about virtual environment

% some new papers from 2024?
\chapter{Gesture Based Robot Control}

\section{Overview}
Two gestures are required to successfully perform the "Pick And Place" task:
a deictic gesture to indicate the position and another gesture for confirmation.\par
During the selection of the object and the target location, the user should face the camera, with all objects lying on the floor between him and the camera. The scene is displayed on the rViz screen. Once the objects are detected by the vision system, their images are marked in blue.\par
The user can then initiate the selection of an object by pointing at it with his right hand. To confirm the gesture, he raises his left hand while still pointing at the object.\par
The object closest to the intersection of the pointing ray and the floor is selected. Its image is marked in red.\par
We can then specify the target location of the selected object. The user determines the location in the same way as before: by pointing to the location and raising his hand.\par
The target location is on the floor and  has to be selected inside the safety frame that is shown in the rViz. The frame represents a space that is safe for the robot to move around, there are no obstacles except for the detected objects.\par
Once the target location is selected, it is marked in red in rViz, gesture detection is completed and the resulting data is sent to the robot.\par

\section{Gestures}

\subsection{Pointing Gesture}
The user can choose from three types of pointing gestures. Each type is represented by a pair of joints that corresponds to the pointing ray. The first joint in the pair is the origin and the second determines the direction of the ray:\par
\begin{itemize}
	\item Head, wrist (default option).
    \item Elbow, wrist.
    \item Shoulder, hand.
\end{itemize}

The pointing gesture indicates the point where the pointing ray intersects the floor and allows us to select an object or its target location.\par
The pointing ray is displayed from its origin to the intersection with the floor.\par
The pointing gesture has to be performed with the right hand and the first joint has to be positioned higher than the second. These constraints help to reduce the number of falsely detected gestures.\par

\subsection{Raising Hand Gesture}
The hand gesture consists of lifting the hand. It has to be performed with the left arm and the hand  has to be raised above the head.\par
When pointing with the right arm, the user confirms the pointing gesture by raising the left hand. If no pointing gesture is performed in the moment, the raising hand gesture is ignored.\par

\section{Gesture Recognition with ORBBEC Astra SDK}

\subsection{Limitations}
ORBBEC Astra SDK provides tools for skeleton recognition and person tracking. The maximum distance for skeleton recognition is 4 meters. Multiple persons can be tracked at the same time.\par
The skeleton is represented by a set of joints and their positions. The head corresponds to single joint, eye positions and other details are unavailable.\par
Three joints are given for each arm: shoulder, elbow and hand. The fingers are not recognizable.\par
SDK also supplies the detection of the grip gesture. I considered using it as a confirmation gesture but preferred the hand raising gesture because the grip was often not detected.\par

\subsection{Occlusion}
The most common cause of gesture recognition errors is occlusion. Only relatively small objects (not taller than 20 cm) were used for the experiments. If the objects were larger, many false detections occurred because the objects often obscured parts of the person's body and the joints of the corresponding skeleton were not correctly identified.\par
Occlusion also often occurs when more than one person is in front of the camera. Therefore, it is better to perform the Pick And Place task when only one person is in the scene because it makes gesture recognition more robust.\par
However, if more than one person is present, each person can perform a pointing gesture to select an object or a target location.\par

\subsection{Recommendations}
When I tested skeleton recognition with a single person and small objects, most of the errors were caused by the person's posture.\par
For best results, the person should be facing the camera, not turning the body or crossing his limbs, as this can lead to errors such as interrupted tracking of the person, misidentification of joints and false recognition of gestures.\par
In case of difficulties with skeleton recognition, it may help if the person moves closer to the camera or extends his arms out to the sides.\par
\chapter{Design of Robotic System}

\section{Basic Structure}

The proposed robotic system contains two main components: a vision system and a mobile robotic manipulator.
The vision system is static, camera remains at the designated location during the task.\par
The robot starts at the initial position where it waits for messages from the vision system. Once the object and target position are selected and a result message is received, the robot navigates to the object, moves it to the target position and returns to the initial position.\par

\subsection{Vision System}

The main purpose of the vision system is to interpret the environment: to detect objects and the skeleton of a person with its gestures. The position of the objects and the tracked person is limited - outside a given frame, the detection is unreliable.\par
For the ORBBEC Astra, the distance must be less than 4 meters. Therefore, I decided to use a vision system separate from the mobile robotic manipulator. Otherwise, with a camera attached to the robot, only objects very close to it could be detected.\par
The depth camera is connected to a laptop where the input data from the camera is processed. The resulting message is sent to the robot's desktop computer.\par

\subsection{Mobile Robot Manipulator}
The robot consists of a mobile vehicle with a robotic arm and a gripper.\par
The vehicle is equipped with a laser scanner that enables localization and safe autonomous navigation in the environment. The on-board computer serves as the robot's ROS master and is connected to a desktop computer.\par
The arm with gripper is attached to the vehicle. It needs to be set up so that there are no collisions with the robot or the floor during manipulation. The reach of the arm should also be limited so that it does not move within the field of view of the scanner, as this would trigger an emergency stop.\par
The arm computer controls both the arm and the gripper and is available via ROS and Dashboard Server.\par


\section{Hardware}

\subsection{Neobotix MP-500}
The MP-500 mobile robot is a differential-wheeled robot with two large drive wheels and one small one at the rear. It is one of the most robust Neobotix mobile robots with a weight of 70 kg.\par
Its main components are a mobile platform, laser scanner, on-board computer, battery pack, manual charger and wireless joystick. Additional components can be attached to the mobile platform.\par
The robot can be used for material transport, with a load capacity of 80 kg. It is designed for indoor operation and has a speed of up to 1.5 m/s.\par
A Sick S300 safety laser scanner with a maximum range of 30 meters is mounted in the front of the mobile platform. The scanner provides data that is used for localization, navigation and collision avoidance.\par
Detection of a person or obstacle in the safety field immediately triggers an emergency stop.\par


\subsection{Robotic Arm UR5}
The Universal Robots UR5 manipulator consists of a robotic arm, a control box with a teaching pendant and a battery.\par
The six-axis arm is composed of extruded aluminum tubes and rotational joints (Base, Arm, Elbow, Wrist 1, Wrist 2, Wrist 3). The Base is the first joint of the kinematic chain in which the arm is mounted to a fixed surface or a mobile platform. The last joint to which the tool is attached is Wrist 3.\par
All joints have a motion range of 360 degrees. The reach of the arm is 0,85 m from the center of the base, the area directly above and below the base is out of reach. The weight is 18.4 kg and the maximum payload is 5 kg.\par
The teaching pendant provides a GUI for control of the arm, commands can also be sent remotely using dedicated ports.\par

\subsection{Weiss Robotics GRIPKIT}

A two-finger gripper is connected to the UR5 arm using the Weiss Robotics GRIPKIT module. Its maximal opening stroke is ? TODO \par

\subsection{ORBBEC Astra camera}
TODO \\

\subsection{Computers and network}
For the vision system, an Acer TravelMate P214 notebook is used. The ORBBEC Astra camera is connected via USB.\par
A Lenovo ThinkStation P330 desktop computer controls the mobile manipulator. It is connected to the Neobotix MP-500 mobile robot via an Ethernet cable.\par
The connection between the computers is established via WiFi, messages are sent using SSH.\par

\section{System Software}
TODO\\
Ubuntu 20.04.\\
Robot Operating System, Noetic.\\

\chapter{Implementation}

\section{Vision system}

\subsection{Image Processing with ORBBEC Astra Camera}
The image data is sent from the camera to the connected laptop for processing.\par
There is the ROS Master and several other individual ROS nodes running on the notebook. Some ROS nodes are involved in image processing, while others provide tools such as geometric calculations or displaying detected objects and skeletons on the rViz screen.\par
ROS nodes communicate with each other using ROS messages and share information about the progress of their subtasks.\par
I used two main tools to process the camera data: the ROS Astra camera driver "ros\_astra\_camera" for object detection and the "ORBBEC Astra SDK" for skeleton detection.\par
Both tools use OpenNI as an intermediate layer to access the camera data. I couldn't run them at the same time because it led to runtime errors, so I decided to split the process into two separate phases.
First, the "ros\_astra\_camera" driver is started. Once all objects are detected, the driver stops.\par
In the second phase, data is exposed using the "ORBBEC Astra SDK" until both pointing gestures are confirmed and the result is sent to the robot. Since gesture recognition is performed within the ROS system, additional ROS packages were needed to publish the body tracking data provided by the SDK as ROS messages.\par

\subsection{Installation}

\subsubsection{ROS Noetic}
The notebook with Ubuntu 20.04 was used, for which the recommended version of the ROS distribution is ROS Noetic. I followed the instructions from http://wiki.ros.org/noetic/Installation/Ubuntu to download and install the ROS Noetic package.\par

\subsubsection{Astra and OpenNI SDKs}
For the ORBBEC Astra camera, I needed to install the Astra SDK and the OpenNI SDK for Linux.\par
Both SDKs are available at https://www.orbbec.com/developers.\par

\subsubsection{ROS Driver for Astra camera}
I downloaded the ROS driver package\\
from https://github.com/orbbec/ros\_astra\_camera and installed the dependencies according to the instructions on http://wiki.ros.org/astra\_camera.\par
The "ros\_astra\_camera" package supports the ROS distributions Kinetic and Melodic. I needed to find and test multiple versions of the "ros-*-libuvc-*" libraries, as they were not released specifically for ROS Noetic.\par
This problem was already solved on the ORBBEC GitHub page, so I followed the advice and installed the missing dependencies using:\par

\begin{lstlisting}[language=bash]
  $ apt install ros-noetic-rgbd-launch libuvc-dev
\end{lstlisting}


I built the package with "catkin\_make" command and was able to run code samples that show the camera data on the screen.\par

\subsubsection{ROS Packages for Gesture Based Control}
I downloaded three ROS packages from the Shinsel Robots repository on https://github.com/shinselrobots.\par
The "pcl\_object\_detection" package allows object detection in the camera data provided by the ROS Astra driver using the Point Cloud Library.\par
The "astra\_body\_tracker" and "body\_tracker\_msgs" packages publish body tracking data from the Astra SDK as ROS messages.\par
Several environment variables have to be set to indicate the paths to the AstraSDK and OpenNI subfolders.\par
For example, if "/home/user/AstraSDK" is the folder containing the Astra SDK and "/home/user/OpenNI-Linux-x64-2.3.0.66" is the folder containing the OpenNI SDK, the settings can be made by running these commands in the terminal: \par

\begin{lstlisting}[language=bash]
  $ export ASTRA_SDK=/home/user/AstraSDK
  $ export ASTRA_ROOT=/home/user/AstraSDK
  $ export ASTRA_SDK_INCLUDE=/home/user/AstraSDK/include
  $ export ASTRA_SDK_LIB=/home/user/AstraSDK/lib
  $ export OPENNI2_INCLUDE=/home/user/OpenNI-Linux-x64-2.3.0.66
  $ export OPENNI2_REDIST=/home/user/OpenNI-Linux-x64-2.3.0.66/Redist
\end{lstlisting}

\subsection{The Vision System Code}

\subsubsection{Catkin Workspace} 
Catkin is the official build system for ROS. Project packages that are placed in the same catkin workspace can be built all at once.\par
My catkin workspace folder contains following ROS packages: \par

\begin{itemize}
	\item ros\_astra\_camera
    \item task\_execution
    \item rviz\_screen
    \item pcl\_object\_detection
    \item pointing\_gesture
\end{itemize}

\subsubsection{Program Overview} 
  
The main launch file is task\_execution.launch. It starts the ROS Astra driver, rViz and other ROS nodes involved in the task: task\_execution\_node, pcl\_object\_detection\_node and pointing\_gesture\_node.\par
The task\_execution\_node subscribes to ROS messages\\
"pcl\_object\_detection/detected\_objects" and "body\_tracker/intersection".\\
The "pcl\_object\_detection/detected\_objects" message contains an array of coordinates of the detected objects, the "body\_tracker/intersection" message contains the coordinates of the intersection of the pointing ray and the floor.\par
When the intersection message is received for the first time, the nearest object is calculated. The object is represented by its index in the detected object array, which is then published in the "task\_execution/pointed\_object\_index" message.\par
The second received intersection message indicates the target location.
Once received, the node creates a result file and writes the coordinates of all detected objects, the coordinates of the target location and the index of the selected object.\par
Then the node connects to the robot's computer using SSH, transfers the file there and remotely starts the robot's main program.\par

\subsubsection{Object Detection with Point Cloud Library}
Point Cloud Library (PCL) was used as a tool for processing the image data.\par
The pcl\_object\_detection package allows the detection of objects on a flat surface and depends on two PCL packages: pcl\_ros and pcl\_conversions.\par
The pcl\_ros package bridges ROS systems and 3D applications that work with point clouds. It extends the ROS C++ Client Library to support messages with native PCL data types.\par
Conversions between PCL data types and ROS message types are provided by the pcl\_conversions package.\par
The pcl\_object\_detection\_node subscribes to the topic \\"/astra\_camera/depth/points", which is published by the ROS Astra camera driver. The received messages represent a point cloud with no color information.\par
The point cloud is processed using PCL: the data is filtered with VoxelGrid and used for plane segmentation. The remaining data, which is out-of-plane, is divided into clusters.
Each cluster can eventually be considered as a detected object. We can specify parameters that limit the height and width of the detected objects.\par
All clusters that meet the given conditions correspond to detected objects and their properties are published in the "pcl\_object\_detection/detected\_objects" message.\par

\subsubsection{Code Adaptation to Experimental Environment}
Even if the scene with the objects was static, different objects were detected in consecutive depth clouds. The detection reliability decreased with the person moving in the scene.\par
I modified the original pcl\_object\_detection package from Shinsel Robots to improve the quality of object detection in my experimental environment. I also needed to implement a mechanism for switching between the ROS Astra driver and the Astra SDK to avoid problems with camera data accessibility.\par
The object size is limited to 20 cm and the object can be detected only when lying on the floor. Furthermore, a detection frame has been specified as a constraint for the object's location on the floor. An object outside the detection frame is ignored.\par
To ensure that the same objects are detected when the experiment is repeated in the same scene, the total number of objects is fixed and the program is modified:\par
The received point clouds are processed one by one, with some objects detected in each cloud. \par
If the number of detected objects does not equal the total number, the result is ignored and processing continues with the next cloud.\par
Otherwise, the number of detected objects is correct, which usually means that the correct objects have been detected. Their data is published in the "pcl\_object\_detection/detected\_objects" message. In addition, the message "object\_detection\_done" is published, which indicates that the ROS Astra driver is no longer needed.\par 
I added a subscriber to the ROS Astra driver to the "object\_detection\_done" message. Once the message is received, the driver is shut down.\par
The package has also been extended with code to display detected objects in rViz, mark them with a color, and change the color of the object that was selected by the pointing gesture.\par

\subsection{Gesture detection}
Packages: 
astra\_body\_tracker:\\
https://github.com/shinselrobots/astra\_body\_tracker\\
Publisher of ROS topic for body tracking information (from the ORBBEC SDK).\\

pointing\_gesture:\\
modified astra\_body\_tracker package to get skeleton data;\\
added code to detect gestures, rViz markers, ....\\

Publisher of pointing\_gesture topic (as geometry\_msgs).\\


\subsection{Object selection and target location}
task\_control\_node (will be renamed):\\
Subscriber to object detection and pointing gesture topics;\\
provides calculations of pointing ray intersection and selection of object.\\
Sends data to mobile manipulator PC over SSH (coordinates of objects and target location, info about selected object).\\

\subsubsection{TODO - some notes}
Why I choose ORBBEC Astra camera over Kinect ONE (v2):\\
difficult installation of tools and libraries for a ROS Interface to the Kinect One (dependencies on ROS Hydro/Indigo distribution, no available packages for ROS Noetic).\\
new wrapper!!!!!!!!!
Installation of ORBBEC SDK for Linux and dependencies (OpenNI2, libsfml-dev, ...).\\

 Package ros\_astra\_camera:\\
 https://github.com/orbbec/ros\_astra\_camera\\
 OpenNI2 ROS wrapper for Orbbec 3D cameras.\\



\section{Navigation of Autonomous Vehicle}

\subsection{Installation}
Neobotix:

Packages:
https://github.com/neobotix/

ros-noetic-amcl, ros-noetic-map-server, ros-noetic-move-base, ...

\subsection{Map of Environment}
Mapping procedure, selecting the map for navigation, visualization with RViz...\\

\subsection{Navigation to Goal}
Goal definition, movement (path, obstacle avoidance, ...).

\section{Object Manipulation}


\subsection{Installation}
Universal Robots:\\

Packages:\\
Universal\_Robots\_ROS\_Driver
https://github.com/UniversalRobots/Universal\_Robots\_ROS\_Driver\\

Universal\_Robots\_Client\_Library\\
https://github.com/UniversalRobots/Universal\_Robots\_Client\_Library\\

ur5\_moveit\_config\\
https://github.com/ros-industrial/universal\_robot/tree/noetic-devel/ur5\_moveit\_config\\



\subsection{Mobile Manipulator URDF}
URDF for Neobotix, UR5 and gripper.\\

\subsection{MoveIt Setup Assistant}
How to create config and set up arm positions.\\
How to set up arm limits.\\
Simulation in rViz.\\

\subsection{Code}
ur\_robot\_driver;\\
ROS.urp;\\
move\_it\_planning;\\
trajectory commands;\\

\subsection{Objects coordinates}
approximate coordinates of objects obtained from the vision system; \\
robot navigates to objects; \\
exact objects coordinates from LIDAR (lidar\_scan topic subscriber).\\











\chapter{Experiments}

Experiments descriptions:

Experiments with different ways of using deictic gestures:\\
- a pointing ray calculated from a pair of skeleton coordinates (head - hand, elbow - wrist, shoulder - wrist)\\
- pointing with or without visual feedback (pointed ray shown in rViz)\\

Experiments measurements:\\
...\\

result evaluation; what went wrong; future work, possible improvements 
\chapter*{Conclusion}
\addcontentsline{toc}{chapter}{Conclusion}

\chapter{Appendix}

%%% Bibliography
%%% Bibliography (literature used as a source)
%%%
%%% We employ biblatex to construct the bibliography. It processes
%%% citations in the text (e.g., the \cite{...} macro) and looks up
%%% relevant entries in the bibliography.bib file.
%%%
%%% See also biblatex settings in thesis.tex.

%%% Generate the bibliography. Beware that if you cited no works,
%%% the empty list will be omitted completely.

% We let bibliography items stick out of the right margin a little
\def\bibfont{\hfuzz=2pt}

%\printbibliography[heading=bibintoc]


%%%%

%\begin{document}

%\maketitle

\section{First section}


\medskip


\printbibliography

%%%%
%%% If case you prefer to write the bibliography manually (without biblatex),
%%% you can use the following. Please follow the ISO 690 standard and
%%% citation conventions of your field of research.

% \begin{thebibliography}{99}
%
% \bibitem{lamport94}
%   {\sc Lamport,} Leslie.
%   \emph{\LaTeX: A Document Preparation System}.
%   2nd edition.
%   Massachusetts: Addison Wesley, 1994.
%   ISBN 0-201-52983-1.
%
% \end{thebibliography}


%%% Figures used in the thesis (consider if this is needed)
\listoffigures

%%% Tables used in the thesis (consider if this is needed)
%%% In mathematical theses, it could be better to move the list of tables to the beginning of the thesis.
\listoftables

%%% Abbreviations used in the thesis, if any, including their explanation
%%% In mathematical theses, it could be better to move the list of abbreviations to the beginning of the thesis.
\chapwithtoc{List of Abbreviations}

%%% Doctoral theses must contain a list of author's publications
\ifx\ThesisType\TypePhD
\chapwithtoc{List of Publications}
\fi

%%% Attachments to the thesis, if any. Each attachment must be referred to
%%% at least once from the text of the thesis. Attachments are numbered.
%%%
%%% The printed version should preferably contain attachments, which can be
%%% read (additional tables and charts, supplementary text, examples of
%%% program output, etc.). The electronic version is more suited for attachments
%%% which will likely be used in an electronic form rather than read (program
%%% source code, data files, interactive charts, etc.). Electronic attachments
%%% should be uploaded to SIS. Allowed file formats are specified in provision
%%% of the rector no. 72/2017. Exceptions can be approved by faculty's coordinator.
\appendix
\chapter{Attachments}

\section{First Attachment}

\end{document}
